\documentclass[11pt,a4paper]{article}
\usepackage[T1]{fontenc}
\usepackage[utf8]{inputenc}
\usepackage{lmodern}
\usepackage{amsmath,amssymb,amsthm,mathtools,mathrsfs}
\usepackage{graphicx,float,xcolor,multicol}
\usepackage[shortlabels]{enumitem}
\usepackage[margin=27mm]{geometry}
\usepackage{microtype,needspace,placeins}
\usepackage{tikz,tikz-cd}
\usetikzlibrary{arrows.meta,calc,positioning,patterns,decorations.pathreplacing}
\usepackage[colorlinks=true,linkcolor=teal,urlcolor=teal,citecolor=teal]{hyperref}
\setlength{\parindent}{0pt}
\setlength{\parskip}{.6em}
\setcounter{secnumdepth}{0}
\allowdisplaybreaks
\raggedbottom
\providecommand{\cross}{\times}
\DeclareUnicodeCharacter{2020}{\ensuremath{\dagger}}
\DeclareMathOperator{\supp}{supp}
\DeclareMathOperator*{\esssup}{ess\,sup}
\providecommand{\chapter}{\section}
\newenvironment{tool}[2]{\subsection*{Tool #1 --- #2}}{}
\newcommand{\ToolRef}[1]{Tool #1}
\newcommand{\FigureTag}[2]{\textit{Figure #1}}
\newcommand{\Result}[1]{\subsection*{#1}}
\newcommand{\Application}[1]{\subsubsection*{Application\if\relax\detokenize{#1}\relax\else\space--- #1\fi}}
\newcommand{\ProofHeading}{\par\textit{Proof.}\space}
\newcommand{\ProofOf}[1]{\par\textit{Proof of #1.}\space}
\newcommand{\RemarkHeading}{\par\textbf{Remark.}\space}
\newcommand{\NoteHeading}{\par\textbf{Note.}\space}
\newcommand{\ExerciseHeading}{\par\textbf{Exercise.}\space}

% Rendering support only; mathematical text is in each independent post.
\usepackage{booktabs,array,longtable}
\definecolor{HandoutBlue}{HTML}{2F6690}
\definecolor{HandoutNavy}{HTML}{17365D}
\newtheorem{theorem}{Theorem}
\newtheorem{lemma}[theorem]{Lemma}
\newtheorem{proposition}[theorem]{Proposition}
\newtheorem{corollary}[theorem]{Corollary}
\newtheorem{conjecture}[theorem]{Conjecture}
\newtheorem{definition}{Definition}
\newtheorem{example}{Example}
\newtheorem{namedtoolcounter}{Tool}
\newenvironment{namedtool}[1]{\begin{namedtoolcounter}[#1]}{\end{namedtoolcounter}}
\newtheorem*{claim}{Claim}
\newtheorem*{remark}{Remark}
\newtheorem*{exercise}{Exercise}
\newtheorem*{question}{Question}
\newtheorem*{observation}{Observation}
\newcommand{\SourceChapter}[2]{\section*{#2}}
\newcommand{\Lecture}[2]{\section*{Lecture #1: #2}}
\newcommand{\unclear}[1]{\textit{[unclear: #1]}}
\newcommand{\sourcegap}{\textit{[unfinished in the source]}}
\DeclareMathOperator{\dist}{dist}
\DeclareMathOperator{\id}{id}
\DeclareMathOperator{\Hol}{Hol}
\DeclareMathOperator{\Per}{Per}
\DeclareMathOperator{\Fix}{Fix}
\DeclareMathOperator{\diam}{diam}
\DeclareMathOperator{\Var}{Var}
\DeclareMathOperator{\Leb}{Leb}
\DeclareMathOperator{\Hom}{Hom}
\DeclareMathOperator{\End}{End}
\DeclareMathOperator{\Ann}{Ann}
\DeclareMathOperator{\Spec}{Spec}
\DeclareMathOperator{\Max}{Max}
\DeclareMathOperator{\Ob}{Ob}
\DeclareMathOperator{\Tor}{Tor}
\DeclareMathOperator{\Ass}{Ass}
\DeclareMathOperator{\Mor}{Mor}
\DeclareMathOperator{\Fun}{Fun}
\DeclareMathOperator{\Nat}{Nat}
\DeclareMathOperator{\Cone}{Cone}
\DeclareMathOperator{\MCG}{MCG}
\DeclareMathOperator{\Aut}{Aut}
\DeclareMathOperator{\Deck}{Deck}
\DeclareMathOperator{\SL}{SL}
\DeclareMathOperator{\PSL}{PSL}
\DeclareMathOperator{\Area}{Area}
\DeclareMathOperator{\im}{im}
\DeclareMathOperator{\PSh}{PSh}
\newcommand{\CRing}{\mathsf{CRings}}
\newcommand{\AMod}{A\text{-}\mathsf{Mod}}
\newcommand{\Ab}{\mathsf{Ab}}
\newcommand{\Z}{\mathbb Z}
\newcommand{\Q}{\mathbb Q}
\newcommand{\R}{\mathbb R}
\newcommand{\C}{\mathbb C}
\newcommand{\N}{\mathbb N}
\newcommand{\kk}{\mathbb k}
\renewcommand{\aa}{\mathfrak a}
\newcommand{\bb}{\mathfrak b}
\newcommand{\pp}{\mathfrak p}
\newcommand{\qq}{\mathfrak q}
\newcommand{\mm}{\mathfrak m}
\newcommand{\nn}{\mathfrak n}
\newcommand{\rad}{\mathrm r}
\newcommand{\cat}[1]{\mathsf{#1}}
\setcounter{secnumdepth}{0}
\allowdisplaybreaks

\graphicspath{{figures/lemma-book/}}
\begin{document}
\section*{Rearrangement, Convexity, and Power Means}
% BLOG-CONTENT-BEGIN
\setcounter{theorem}{54}\begin{lemma}[Popoviciu's theorem]
Let $f:\mathbb R\to\mathbb R$ be convex on an interval $I$, and let $x_1,x_2,x_3\in I$.
Then
\begin{align*}
f(x_1)+f(x_2)+f(x_3)+3f\left(\frac{x_1+x_2+x_3}{3}\right)
\geq{}&2f\left(\frac{x_1+x_2}{2}\right)+2f\left(\frac{x_2+x_3}{2}\right)\\
&+2f\left(\frac{x_3+x_1}{2}\right).
\end{align*}
\end{lemma}
This classic was formulated by the Romanian mathematician Tiberiu Popoviciu.
\paragraph{Application.}
For real $x,y,z$, prove
$|x|+|y|+|z|-|x+y|-|y+z|-|z+x|+|x+y+z|\geq0$.
Let $f(a)=|a|$, convex on $(-\infty,\infty)$.
By Popoviciu's theorem,
\[
|x|+|y|+|z|+3\left|\frac{x+y+z}{3}\right|
\geq2\left|\frac{x+y}{2}\right|+2\left|\frac{y+z}{2}\right|+2\left|\frac{z+x}{2}\right|.
\]
This is exactly the desired inequality.

\subsection{Rearrangement and convexity}
\setcounter{theorem}{84}\begin{lemma}[Rearrangement inequality]\label{lemma:lb1-87}
Consider two collections of real numbers in increasing order,
$a_1\leq a_2\leq\cdots\leq a_n$ and $b_1\leq b_2\leq\cdots\leq b_n$.
For any permutation $(a'_1,\ldots,a'_n)$ of $(a_1,\ldots,a_n)$,
\[a_1b_1+a_2b_2+\cdots+a_nb_n\geq a'_1b_1+a'_2b_2+\cdots+a'_nb_n.\]
\end{lemma}
\paragraph{Application.}
Let $a,b,c$ be positive numbers with $abc=1$. Prove that
\[\frac ab+\frac bc+\frac ca\geq a+b+c.\]
\paragraph{Solution 1.}
Apply the rearrangement inequality to
\[(a_1,a_2,a_3)=\left(\sqrt[3]{\frac ab},\sqrt[3]{\frac bc},\sqrt[3]{\frac ca}\right),
\quad(b_1,b_2,b_3)=\left(\sqrt[3]{\left(\frac ab\right)^2},\sqrt[3]{\left(\frac bc\right)^2},\sqrt[3]{\left(\frac ca\right)^2}\right),\]
and the permutation $(a'_1,b'_1,c'_1)=(\sqrt[3]{b/c},\sqrt[3]{c/a},\sqrt[3]{a/b})$, to derive
\[\frac ab+\frac bc+\frac ca\geq\sqrt[3]{\frac{a^2}{bc}}
+\sqrt[3]{\frac{b^2}{ca}}+\sqrt[3]{\frac{c^2}{ab}}=a+b+c,\]
since $abc=1$. This completes the proof.
Two other proofs are available; the next uses AM--GM.
\paragraph{Solution 2.}
\[
\begin{aligned}
\frac13\left(\frac ab+\frac ab+\frac bc\right)&\geq\sqrt[3]{\frac{a\cdot a\cdot b}{b\cdot b\cdot c}}=\sqrt[3]{\frac{a^2}{bc}}=a,\\
\frac13\left(\frac bc+\frac bc+\frac ca\right)&\geq\sqrt[3]{\frac{b\cdot b\cdot c}{c\cdot c\cdot a}}=\sqrt[3]{\frac{b^2}{ca}}=b,\\
\frac13\left(\frac ca+\frac ca+\frac ab\right)&\geq\sqrt[3]{\frac{c\cdot c\cdot a}{a\cdot a\cdot b}}=\sqrt[3]{\frac{c^2}{ab}}=c.
\end{aligned}
\]
Add all three to obtain $a/b+b/c+c/a\geq a+b+c$.
The third solution follows from Muirhead's theorem.
% Source page 46 - right
\paragraph{Solution 3.}
\[\frac ab+\frac bc+\frac ca=a^2c+b^2a+c^2b
\mathrel{\stackrel{?}{\geq}}(abc)^{2/3}(a+b+c)
=\sum_{\mathrm{cyc}}a^{5/3}b^{2/3}c^{2/3}.\]
Since $(2,1,0)\succ(5/3,2/3,2/3)$, the result follows. This completes the proof.
% Author check: the source applies Muirhead to a three-term cyclic sum; its usual symmetric-sum hypothesis is not addressed.
\paragraph{A two-line proof of the AM--GM inequality.}
We have $x_i^{t_i}=e^{t_i\log x_i}$. Since $e^x$ is convex, we can deduce
\[x_1^{t_1}x_2^{t_2}\cdots x_n^{t_n}
=e^{t_1\log x_1+\cdots+t_n\log x_n}
\leq t_1e^{\log x_1}+\cdots+t_ne^{\log x_n}=t_1x_1+\cdots+t_nx_n.\]
Now put $t_1=t_2=\cdots=t_n=1/n$ to obtain, by Jensen's inequality,
\[\frac{x_1+x_2+\cdots+x_n}{n}\geq\sqrt[n]{x_1x_2\cdots x_n}.\]
This completes the proof.
% Author check: the initial general-weight formula does not state that the weights are nonnegative and sum to one.
\paragraph{Challenge.}
Define $x_k=(a^k+b^k+c^k)/3$ for $a,b,c>0$. Prove the following inequality:
\[\frac13\left(\frac{x_1^5}{x_2x_3}+\frac{x_2^2}{x_1x_3}+\frac{x_3}{x_1x_2}\right)\geq1.\]
The proof takes only two lines.
% Author check: the strict sign in the source is retained; a=b=c gives equality.
Convexity and concavity can make apparently difficult inequalities collapse quickly.

% Source page 49 - left
Here is a two-step solution to the challenge above.
\[
\begin{aligned}
x_3^2\geq x_2^3\geq x_1^6
&\quad\Longrightarrow\quad
\frac13(x_3^2+x_2^3+x_1^6)\left(\frac12+\frac13+\frac16\right)
\geq\frac{x_3^2}{2}+\frac{x_2^3}{3}+\frac{x_1^6}{6}\geq x_1x_2x_3,\\
\frac{x_3^2+x_2^3+x_1^6}{x_1x_2x_3}\geq3
&\quad\Longrightarrow\quad
\frac{x_1^5}{x_2x_3}+\frac{x_2^2}{x_1x_3}+\frac{x_3}{x_1x_2}\geq3.
\end{aligned}
\]
This completes the proof in two steps.
% Author check: the source uses a strict > in the weighted-sum comparison; the direction and strictness are preserved for review.

For $a,b,c>0$ with $abc=1$, prove $\sqrt{a+3}+\sqrt{b+3}+\sqrt{c+3}\geq6$.
Since $ab+bc+ca\geq3$ by AM--GM,
\[
\sum_{\mathrm{cyc}}\sqrt{a+3}
\geq3\sqrt[6]{28+3(ab+bc+ca)+9(a+b+c)}
\geq3\sqrt[6]{28+9+27}=6.
\]
% Author check: source uses sixth roots in this displayed AM--GM calculation and equality ab+bc+ca=3; retained.

For positive $x,y,z$, prove
$\sqrt{x+y+z}\left(\sum_{\mathrm{cyc}}\dfrac{\sqrt y}{z+x}\right)\geq3\sqrt3/2$.
Normalize $x+y+z=1$, and set $f:(0,1)\to\mathbb R$, $f(x)=\sqrt x/(1-x)$.
Since $f$ is convex, $f(x)+f(y)+f(z)\geq3f(1/3)$, which gives the result.
% Author check: asserted convexity on the entire interval is preserved.
% BLOG-CONTENT-END
\end{document}
